\section{系统与威胁模型}\label{sec:overview}

\subsection{系统边界与参与者}\label{sec:overview:layout}

Constitutional ARI-RQGM 治理围绕自主研究 executor 的制度。executor 在
checkpoint 中产生类型化研究产物、测量、review record 与 manuscript
candidate。治理系统观察这些记录,并控制哪些 prompt、component 与
utility-policy version 可以服务。host OS、model weight 以及能够修改程序的
人类不在控制范围内。

参与者分为四类。\emph{research actor} 负责提议、评估、辩护与写作;
\emph{judiciary actor} 收集证据、起诉、为 governance motion 辩护并裁决;
\emph{meta actor} 生成 candidate prompt 和 policy,或选择 replay case;
\emph{fixed actor} 验证宪制规则、应用 registry 转换、修复受影响记录并
核验最终主张。参与者只能在声明的角色内生成带 hash 的记录。当前成果物
没有实现数字签名或角色密钥管理。产生建议的权限并不意味着执行建议的
权限。

\subsection{威胁模型}\label{sec:overview:pipeline}

首要威胁是递归制度优化:可演化角色可能寻求让自身输出看起来更好的改变,
而非提高研究可靠性。具体威胁包括 score inflation、metric gaming、
evidence omission、prompt injection、proposer 与 evaluator 共谋、meta agent
请求比 incumbent 更宽的工具权限、judge 裁决针对自己的起诉,以及退役
prompt 通过派生 score 继续施加影响。

模型也覆盖意外故障。进程可能在 registry event 之间停止,model 可能输出
错误格式,anchor corpus 可能缺失,prompt candidate 可能违反 output schema,
保存的证据也可能缺少重新评分所需的 raw axis。这些故障均不得形成只应用
一半的制度。

恶意 host administrator、对固定程序的任意修改、底层 hash function 被
攻破、保存点与本地固定值一起回退到相同旧前缀、provider 秘密修改 model
weight,以及不可逆的外部 API、设备或生产数据操作不在威胁模型内。改变
宪制程序被视为人类修正,而不是运行中系统能够阻止的攻击。

\subsection{权限与验证边界}\label{sec:overview:verification}

权限单向流动。带模型的角色可以创建 proposal、attack、defense、judgment
与 candidate,但确定性内核不会把它们直接变成 registry state。内核检查
record、acting role、请求 edge、epoch timing、capability subset、
provenance 与 audit continuity。transition engine 是唯一 registry writer,
其输出在 commit 前还会再次接受检查。

两条不对称故障规则在保留有用研究的同时阻止不确定制度变更。研究路径上的
governance hook 为 \emph{fail-open}:故障被记录,只有无外部副作用沙箱内的
候选生成和本地计算可以继续。改变状态的 governance operation 为
\emph{fail-closed}:无法证明证据或合法性时,incumbent 制度继续服务。
不可逆外部操作需要另一个故障时关闭的固定闸门;当前内核不强制该环境条件。

最终 manuscript verifier 属于固定边界。稿件前向声明数值主张标识、公式、
操作数位置、单位和容差;验证器依据执行证据重算这个封闭公式集合。它不从
自然语言推断任意主张,也不证明非数值主张或执行记录本身的真实性。writer
与 reviewer 不能改变 verifier、输入或 verdict。

\subsection{状态、来源与重放}\label{sec:overview:checkpoint}

checkpoint 是制度 identity 的单位。append-only event 是权威记录;registry、
epoch 与 erasure snapshot 是可以通过 replay 重建的派生 view。每条 governed
record 带有 epoch、serving component、prompt hash 与相关 policy hash。
audit event 构成 predecessor-linked hash chain。因此读者能够区分发生了
什么、由哪个制度产生,以及后续转换如何改变其 eligibility。本地链能检测
中间删除或修改;若本地边界固定值仍存在,也能检测更短日志。但日志和固定
值一起回退到相同旧的合法前缀时无法检测。当前没有外部签名日志头、见证者
保存点或单调计数器,哈希也不是作者签名。

epoch transaction 使用 prepare、close、open 与 commit event。replay 忽略
不完整后缀并确定性重做边界。founding registration 同样 write-once:
已 commit 的 transaction 只被重放,不会重复执行。logical erasure 不删除
字节,而是增加 staleness 与 eligibility 字段,从而保留审计退役本身所需的
证据。

纪元身份先拆分再合成。政策指纹绑定服务制度以及包含阈值与预算在内的全部
已解析治理设置；执行指纹绑定声明的模型、解码、工具、环境与数据快照。
缺少提供者或环境修订时会显式标为未解析。前者支持制度比较，后者暴露同一
政策是否运行在不同的声明执行基底上；二者都不能把不透明模型服务变成可
复现二进制。

\subsection{设计原则}\label{sec:overview:principles}

架构遵循五项原则。

\begin{enumerate}
  \item \textbf{固定法律,演化制度。}宪制、转换 topology 与最终
        verification path 位于 runtime-mutable artifact 集合之外。
  \item \textbf{冻结的服务纪元。}serving version 只在边界改变。T16
        紧急隔离会立即形成边界并开启新纪元指纹，前后节点具有不同制度标识。
  \item \textbf{建议与权限分离。}proposer、judge 与 meta agent 输出记录;
        只有经内核验证的 transition engine 可以写 registry。
  \item \textbf{保留历史但不继承影响。}退役不删除证据,但受影响 descendant
        必须在影响 future selection 前变为 stale 或被重新评分。
  \item \textbf{主张与证据成比例。}结构保证按保证陈述;机制演示不被包装为
        研究质量提升的证据。
\end{enumerate}
